\documentclass[11pt]{article}
\usepackage{ntex-notes}

\title{Specimen — Linear Algebra, Lecture 7}
\author{Your Name}
\date{1 October 2026}

\begin{document}
\maketitle

\section{Spectral bits}

Let $A \in \mathbb{R}^{n \times n}$ be symmetric with eigenvalues
$\lambda_1 \le \dots \le \lambda_n$ and orthonormal eigenvectors $v_1, \dots, v_n$.
Body text is set in Charter; the notes theme only colours structure.\footnote{A footnote, to check its size and colour.}

\begin{definition}[Rayleigh quotient]
\label{def:rayleigh}
For $x \neq 0$, the \emph{Rayleigh quotient} of $A$ is
\[
  R_A(x) = \frac{x^{T} A x}{x^{T} x}.
\]
\end{definition}

\begin{theorem}[Courant--Fischer, easy half]
\label{thm:cf}
For every $x \ne 0$, $\lambda_1 \le R_A(x) \le \lambda_n$, with equality at
$v_1$ and $v_n$ respectively.
\end{theorem}

\begin{lemma}
\label{lem:expand}
If $x = \sum_{i=1}^{n} c_i v_i$, then $\lVert x \rVert^{2} = \sum_{i} c_i^{2}$ and
$x^{T} A x = \sum_{i} \lambda_i c_i^{2}$.
\end{lemma}

\begin{proof}
By orthonormality $\langle v_i, v_j \rangle = \delta_{ij}$, so
\begin{align}
  x^{T} A x &= \sum_{i,j} c_i c_j \, v_i^{T} A v_j
             = \sum_{i,j} c_i c_j \lambda_j \delta_{ij} \label{eq:expand} \\
            &= \sum_{i=1}^{n} \lambda_i c_i^{2}. \qedhere
\end{align}
\end{proof}

\begin{proposition}
The maximum of $R_A$ over the unit sphere is $\lambda_n$; see \cref{thm:cf}, \cref{lem:expand} and \eqref{eq:expand}, with $R_A$ from \cref{def:rayleigh}.
\end{proposition}

\begin{corollary}
If $A \succeq 0$ then every eigenvalue is non-negative.
\end{corollary}

\begin{example}
For
\[
  A = \begin{pmatrix} 2 & -1 \\ -1 & 2 \end{pmatrix}, \qquad
  \det(A - \lambda I) = (2-\lambda)^{2} - 1 = 0 \implies \lambda \in \{1, 3\}.
\]
\end{example}

\begin{remark}
Nothing here needs $A$ to be real: the same holds for Hermitian matrices with
$x^{*}$ in place of $x^{T}$.
\end{remark}

\subsection{Lists and a table}

What to check when tuning:
\begin{itemize}
  \item colour of bullets and numbers;
  \item spacing around boxes;
  \item how a box looks when it breaks across a page.
\end{itemize}

\begin{enumerate}
  \item Compute the characteristic polynomial.
  \item Find the roots.
  \item Normalise the eigenvectors.
\end{enumerate}

\begin{table}[h]
  \centering
  \caption{Eigenvalues of small test matrices.}
  \begin{tabular}{lcc}
    \toprule
    Matrix & $\lambda_{\min}$ & $\lambda_{\max}$ \\
    \midrule
    $\begin{psmallmatrix} 2 & -1 \\ -1 & 2 \end{psmallmatrix}$ & 1 & 3 \\
    $I_2$ & 1 & 1 \\
    \bottomrule
  \end{tabular}
\end{table}

\end{document}
